% Generated draft: requires mathematical and source review.
\documentclass[11pt]{article}
\usepackage[margin=25mm]{geometry}
\usepackage{fontspec}
\setmainfont{Latin Modern Roman}
\usepackage{amsmath,amssymb,mathtools,unicode-math}
\setmathfont{Latin Modern Math}
\usepackage{enumitem,needspace,xurl,hyperref}
\setlength{\parindent}{0pt}
\setlength{\parskip}{5pt}
\setlength{\emergencystretch}{3em}
\newcommand{\N}{\mathbb{N}}
\newcommand{\Z}{\mathbb{Z}}
\newcommand{\Q}{\mathbb{Q}}
\newcommand{\R}{\mathbb{R}}
\newcommand{\C}{\mathbb{C}}
\newcommand{\F}{\mathbb{F}}
\newcommand{\A}{\mathbb{A}}
\newcommand{\PP}{\mathbb{P}}
\newcommand{\E}{\mathbb{E}}
\newcommand{\eps}{\varepsilon}
\newcommand{\dd}{\,\mathrm{d}}
\begin{document}
\section*{UnsolvedMath ID: 11000098}
\textbf{Source problem code:} catalogue entry
\subsection{Definitions and mathematical statement}
Let \(\Gamma\) be an irreducible arithmetic group of rank at least \(2\), in the source's higher-rank lattice setting, and let \(\operatorname{Mod}_S\) be the mapping class group of the intended surface \(S\). Surface type, marking/boundary conventions, and the precise ambient rank convention are not specified. The rank must be checked in that setting rather than silently interpreted as \(\Q\)-rank; the discussion includes cocompact lattices in \(SU(p,q)\) among its motivating cases.

\emph{Conjecture I, reported in the source as completely proved.} Every homomorphism \(\rho:\Gamma\to\operatorname{Mod}_S\) has finite image:
\[
 \forall\rho\in\operatorname{Hom}(\Gamma,\operatorname{Mod}_S),
 \qquad |\rho(\Gamma)|<\infty.
\]
No injectivity or other representation hypothesis is imposed. The initial extracted ``f'' is read as the missing beginning of ``If.''

The source says elementary combinations of arithmetic-group and mapping-class-group facts prove many cases, but notes cocompact \(SU(p,q)\) lattices as a case where that approach appeared to fail, crediting G. Prasad for the example. It then expressly says the full conjecture has been proved. It credits V. A. Kaimanovich and H. Masur [KaM] with showing that nonelementary mapping-class subgroups are not isomorphic to irreducible arithmetic groups of rank at least \(2\). Here nonelementary is described as containing two pseudo-Anosov elements with disjoint fixed-point sets in the Thurston boundary. The proof uses random walks and H. Furstenberg's arithmetic-group results; combined with Margulis's finiteness theorem and Ivanov's techniques, the source says it proves the conjecture. It also credits B. Farb and H. Masur [FM] with a different proof.

The source further credits S-K. Yeung [Y] with finite image for homomorphisms from lattices in \(\operatorname{Sp}(m,1)\) or \(F_4^{-20}\), the latter associated with the Cayley plane. The admissible range of \(m\) and exact hypotheses are absent and must be checked; no unrestricted rank-one extension is inferred. It notes that lattices in remaining simple rank-one Lie groups often have nontrivial homomorphisms to \(\Z\), and hence can yield infinite-image mapping-class homomorphisms in the intended surface setting. ``Often'' is not replaced by ``every.
\subsection{Short English statement}
Every homomorphism from an irreducible higher-rank arithmetic group to a mapping class group should have finite image. The source explicitly reports this conjecture completely proved and gives several proof attributions; an editorial summary incorrectly labels it still open.
\subsection{Sources}
\begin{itemize}
\item Local UnsolvedMath registry, permanent ID 11000098, source code catalogue entry, Conjecture I, PDF page 83. No individual proposer field is supplied.
\item Farb, \emph{Problems on Mapping Class Groups and Related Topics}. Supplied URL: \url{https://www.math.uchicago.edu/~farb/papers/mcgbook.pdf}; recorded access 2026-07-29; extraction pdf-font-structure. The edited volume contains 329 explicitly typeset headings across chapters, with repeated numbering disambiguated by page.
\item V. A. Kaimanovich and H. Masur [KaM]; B. Farb and H. Masur [FM]; S-K. Yeung [Y]. Full entries are missing. The source credits G. Prasad for the cocompact \(SU(p,q)\) example and mentions H. Furstenberg and Margulis in its proof discussion.
\item N. V. Ivanov, \emph{Subgroups of Teichm\"uller Modular Groups}, Translations of Mathematical Monographs, Vol. 115, American Math. Soc., 1992, is explicitly cited for techniques completing the argument.
\item Supplied status solved and research\_classification SOLVED-IN-LITERATURE agree with the source's statement that the conjecture is completely proved. The research summary instead says OPEN-TRIAGE, partial rigidity, and no complete proof through 2026; this directly conflicts with the supplied source and metadata and is not adopted. Generic background still requires a current primary-literature status audit and redistribution review, both NEEDS\_REVIEW. The registry assigns L4 whereas background records default L3.
\end{itemize}
\end{document}
